% Vina Doctor AI Scribe — User Guide
% ===================================
% A comprehensive guide for doctors using Vina Doctor AI Scribe
% to transcribe and structure medical consultation audio into SOAP reports.

\documentclass[12pt,a4paper]{article}

% --- Packages ---------------------------------------------------------------
\usepackage{iftex}
\usepackage{fontspec}
\usepackage{microtype}
\usepackage{geometry}
\usepackage{titlesec}
\usepackage{enumitem}
\usepackage{fancyhdr}
\usepackage{graphicx}
\usepackage{hyperref}
\usepackage{booktabs}
\usepackage{array}
\usepackage{metalogo}

% --- Colors -----------------------------------------------------------------
\usepackage{xcolor}
\definecolor{primary}{HTML}{00595c}
\definecolor{primary-container}{HTML}{0d7377}
\definecolor{secondary}{HTML}{515f74}
\definecolor{tertiary}{HTML}{714600}
\definecolor{error}{HTML}{ba1a1a}
\definecolor{surface}{HTML}{f5fafb}
\definecolor{ondark}{HTML}{ffffff}

% --- Geometry ----------------------------------------------------------------
\geometry{margin=2.5cm}

% --- Fonts -------------------------------------------------------------------
\usepackage{fontspec}
\setmainfont{Roboto}
\setsansfont{Roboto}

% --- Headers / Footers ------------------------------------------------------
\pagestyle{fancy}
\fancyhead[L]{\color{primary}\textbf{Vina Doctor AI Scribe}}
\fancyhead[R]{\thepage}
\fancyfoot[C]{}
\fancyfoot[R]{\color{secondary}\small User Guide — v1.0}
\pagestyle{fancy}

% --- Title formatting -------------------------------------------------------
\titleformat{\section}{\color{primary}\Large\bfseries}{}{0em}{#1}
\titleformat{\subsection}{\color{primary}\large\bfseries}{}{0em}{#1}
\titleformat{\subsubsection}{\color{primary}\normalsize\bfseries}{}{0em}{#1}

% --- Custom commands ---------------------------------------------------------
\newcommand{\cmd}[1]{\texttt{#1}}
\newcommand{\img}[3]{\textit{#3}\includegraphics[width=#1]{#2}}

% --- Hyperref -----------------------------------------------------------
\hypersetup{
  colorlinks=true,
  linkcolor=primary,
  urlcolor=primary-container,
  pdftitle={Vina Doctor AI Scribe — User Guide},
  pdfauthor={Vina Doctor Team},
}

% --- Document ------------------------------------------------------------
\begin{document}

% --- Cover ----------------------------------------------------------------
\begin{titlepage}
	\centering
	\vspace*{2cm}

	% Logo / brand
	\color{primary}\LARGE\bfseries Vina Doctor \\[0.3em]
	\color{primary-container}\large AI Scribe

	\vspace{1.5cm}

	{\color{primary}\Huge\bfseries User Guide}

	\vspace{1cm}
	{\color{secondary}\Large Transforming Consultations into Structured Clinical Notes}

	\vspace{3cm}

	\begin{tabular}{m{2.5cm}m{10cm}}
		\textbf{Version} & 1.0                   \\
		\textbf{Date}    & April 2026            \\
		\textbf{Product} & Vina Doctor AI Scribe \\
	\end{tabular}

	\vspace{2cm}
	\color{secondary}\small This document is confidential and intended for licensed medical professionals.
\end{titlepage}

% --- Table of Contents ------------------------------------------------------
\newpage
\tableofcontents
\newpage

% --- Section 1: Getting Started ---------------------------------------------
\section{Getting Started}

\subsection{What is Vina Doctor AI Scribe?}
Vina Doctor AI Scribe is a web-based portal that uses artificial intelligence to
transform audio recordings of patient consultations into structured clinical notes
using the SOAP format (Subjective, Objective, Assessment, Plan).

\subsection{System Requirements}
\begin{itemize}
	\item A modern web browser: Google Chrome, Mozilla Firefox, Apple Safari, or Microsoft Edge
	\item A stable internet connection
	\item An audio file in one of the following formats: \textbf{MP3, WAV, M4A, OGG, WebM,} or \textbf{FLAC}
	\item A valid account provided by your administrator
\end{itemize}

\subsection{Accessing the Application}
Navigate to your organization's Vina Doctor URL in your browser. You will be
redirected to the login page if you are not already authenticated.

\subsection{Logging In}
\begin{enumerate}
	\item Enter your registered email address in the \textbf{Email} field.
	\item Enter your password in the \textbf{Password} field.
	\item Click \textbf{Sign In}.
	\item On success, you will be taken to the Dashboard.
\end{enumerate}

\subsection{Creating an Account}
If you do not have an account yet:
\begin{enumerate}
	\item Click the \textbf{Register} link below the Sign In button.
	\item Fill in your \textbf{Full Name}, \textbf{Email}, and \textbf{Password} (minimum 8 characters).
	\item Click \textbf{Create Account}.
	\item You will see a success notification and be logged in automatically.
\end{enumerate}

% --- Section 2: Dashboard ----------------------------------------------------
\section{The Dashboard}

After logging in, you land on the \textbf{Dashboard}. This is your home screen.

\subsection{Dashboard Layout}
The Dashboard is divided into three areas:
\begin{itemize}
	\item \textbf{Left sidebar} — navigation menu (hidden on mobile, shown on tablet/desktop).
	\item \textbf{Header row} — a greeting and the \textbf{New Consultation} quick-action button.
	\item \textbf{Main content area} — statistics cards and a list of recent consultations.
\end{itemize}

\subsection{Statistics Cards}
The three cards at the top give you a quick overview:
\begin{tabular}{m{3cm}m{9cm}}
	\textbf{Today's Consultations} & Total number of consultations created today.      \\
	\textbf{Pending Review}        & Consultations still being processed by the AI.    \\
	\textbf{Completed}             & Consultations for which the SOAP report is ready. \\
\end{tabular}

\subsection{Recent Consultations}
Below the statistics, the five most recent consultations are listed as clickable cards.
Click any card to open its detail page. Use the \textbf{View All} link to see the full history.

% --- Section 3: Navigation ----------------------------------------------------
\section{Navigation}

The sidebar provides access to all main sections:

\subsection{Desktop / Tablet}
The sidebar is permanently visible on the left. It contains:
\begin{itemize}
	\item \textbf{Dashboard} — home screen
	\item \textbf{New Consultation} — upload audio and start processing
	\item \textbf{History} — paginated list of all past consultations
	\item \textbf{Settings} — profile, language, and AI model preferences
	\item \textbf{Start Recording} (call-to-action button) — shortcut to New Consultation
	\item \textbf{Log Out} — sign out of your account
\end{itemize}

\subsection{Mobile}
On smaller screens, the sidebar collapses into a hamburger menu
tapped from the top bar. Tapping the hamburger reveals the same navigation links.

% --- Section 4: New Consultation ---------------------------------------------
\section{Creating a New Consultation}

\subsection{Upload Flow}
\begin{enumerate}
	\item Click \textbf{New Consultation} in the sidebar or the \textbf{+} button on the Dashboard.
	\item On the New Consultation page, you will see a drop zone.
	\item \textbf{Drag and drop} your audio file onto the zone, or click \textbf{browse} to select a file from your device.
	\item Once a file is selected, it is displayed as a card showing the filename and size. You can remove it by clicking the $\times$ button and selecting a different file.
	\item Choose an AI model:
	      \begin{itemize}
		      \item \textbf{Optimized Speed} — fast, real-time transcription; suitable for routine visits.
		      \item \textbf{Maximum Accuracy} — deeper clinical analysis; recommended for complex cases.
	      \end{itemize}
	\item Click \textbf{Upload \& Generate Notes}.
	\item The page will redirect to the consultation detail page where you can track progress.
\end{enumerate}

\subsection{Accepted File Formats}
\begin{tabular}{ll}
	\texttt{.mp3}  & MPEG Audio Layer III       \\
	\texttt{.wav}  & Waveform Audio File Format \\
	\texttt{.m4a}  & MPEG-4 Audio               \\
	\texttt{.ogg}  & Ogg Vorbis                 \\
	\texttt{.webm} & WebM Audio                 \\
	\texttt{.flac} & Free Lossless Audio Codec  \\
\end{tabular}

\subsection{AI Model Choices}
\begin{tabular}{m{4cm}m{8cm}}
	\textbf{Optimized Speed}  & Uses the Qwen ASR Flash model. Low latency, real-time-style transcription. Good for follow-up visits and routine checkups.            \\[0.5em]
	\textbf{Maximum Accuracy} & Uses the Qwen3.5-Omni-Flash model. Deep clinical reasoning; best for new patients, complex diagnoses, or multi-symptom presentations. \\
\end{tabular}

\subsection{What Happens After Upload?}
Once you click \textbf{Upload \& Generate Notes}:
\begin{enumerate}
	\item The audio file is saved to the server and a consultation record is created with status \textbf{Pending}.
	\item You are redirected to the consultation detail page.
	\item The AI begins processing the audio in the background.
	\item The detail page automatically refreshes every 3 seconds to show the latest status.
\end{enumerate}

% --- Section 5: Consultation History -----------------------------------------
\section{Consultation History}

\subsection{Viewing Past Consultations}
Click \textbf{History} in the sidebar to see a paginated list of all your consultations.

\subsection{Pagination}
\begin{itemize}
	\item Each page shows up to 20 consultations.
	\item Use \textbf{Previous} and \textbf{Next} buttons to navigate pages.
	\item The page indicator shows the current range, e.g. \textit{Showing 1 to 20 of 156}.
\end{itemize}

\subsection{Table Columns}
\begin{tabular}{m{3cm}m{9cm}}
	\textbf{Date}         & The date and time the consultation was created.                                         \\
	\textbf{Consultation} & A shortened ID (the first 8 characters of the UUID).                                    \\
	\textbf{Status}       & A coloured pill: Pending (blue-grey), Processing (amber), Done (teal), or Failed (red). \\
	\textbf{View}         & A link to open the consultation detail page.                                            \\
\end{tabular}

% --- Section 6: Consultation Detail Page --------------------------------------
\section{Consultation Detail Page}

Click any consultation card (from the Dashboard or History) to open its detail page.

\subsection{Page Layout}
\begin{itemize}
	\item \textbf{Header} — consultation ID and creation date/time.
	\item \textbf{Status badge} — coloured pill reflecting the current state.
	\item \textbf{Progress card} — shown while the AI is working (Pending / Processing).
	\item \textbf{SOAP Report} — shown once processing is complete.
\end{itemize}

\subsection{Status States}

\subsubsection{Pending}
The audio has been received and is waiting in the processing queue.
A spinner card shows: \textit{"AI is queued your recording…"}.

\subsubsection{Processing}
The AI engine is actively transcribing and analysing the audio.
A spinner card shows: \textit{"AI is processing your recording…"}.
This page updates automatically every 3 seconds.

\subsubsection{Done}
The SOAP report is ready and displayed below the header.

\subsubsection{Failed}
Processing could not be completed (for example, due to an API error).
A red error card is shown with a \textbf{Retry} button.
Clicking \textbf{Retry} re-runs the existing audio through the AI — no need to re-upload the file.
The page then returns to the Processing state and polls until complete.

% --- Section 7: SOAP Report -------------------------------------------------
\section{The SOAP Report}

When a consultation reaches the \textbf{Done} status, the SOAP report is displayed on the detail page.

\subsection{Report Structure}
The SOAP report is divided into four sections:

\begin{tabular}{m{3cm}m{9cm}}
	\textbf{Subjective} & What the patient reported — symptoms, duration, medical history, and chief complaint.        \\
	\textbf{Objective}  & Measurable findings — vital signs, physical examination results, and test values.            \\
	\textbf{Assessment} & The clinician's diagnosis or clinical impression, including possible differential diagnoses. \\
	\textbf{Plan}       & Recommended next steps — prescriptions, referrals, follow-up schedule, and lifestyle advice. \\
\end{tabular}

Each section is collapsible. Click the section header to expand or collapse it.

\subsection{Viewing in Your Language}
Above the SOAP sections, three language tabs allow you to switch the report language:
\begin{itemize}
	\item \textbf{Tiếng Việt} — Vietnamese
	\item \textbf{English} — English
	\item \textbf{Français} — French
\end{itemize}
All four SOAP sections and the free-text fields update immediately when you switch tabs.

\subsection{ICD-10 Codes}
If the AI identified diagnostic codes, they appear as teal pills above the SOAP sections,
e.g. \texttt{[A00.1]} \texttt{[B20.5]}. These are clickable (where supported by your
external reference tool).

\subsection{Severity Indicator}
A severity badge appears below the ICD-10 codes. Possible values:
\begin{itemize}
	\item \textbf{Critical} — red badge
	\item \textbf{High} — orange badge
	\item \textbf{Medium} — amber badge
	\item \textbf{Low} — grey badge
\end{itemize}

\subsection{Medications Table}
If the AI generated a prescription, a medications table appears at the bottom of the report:

\begin{tabular}{m{3cm}m{3cm}m{3cm}m{3cm}}
	\textbf{Name} & \textbf{Dosage} & \textbf{Frequency} & \textbf{Duration} \\
	Amoxicillin   & 500 mg          & Three times daily  & 7 days            \\
	Ibuprofen     & 400 mg          & As needed          & 5 days            \\
\end{tabular}

% --- Section 8: Settings ----------------------------------------------------
\section{Settings}

Click \textbf{Settings} in the sidebar to manage your preferences.

\subsection{Doctor Profile}
Display-only fields showing your registered name, specialty, and licence number.
These are set by your administrator.

\subsection{Language Defaults}
Set your preferred languages for dictation input and report output.
Options: Vietnamese, English, French.

\subsection{AI Model Preference}
Choose your default AI model for all new consultations:
\begin{itemize}
	\item \textbf{Optimized Speed} (default)
	\item \textbf{Maximum Accuracy}
\end{itemize}
This preference is saved per-browser using local storage.

\subsection{Updating the DashScope API Key}
If your organisation uses a custom DashScope API key:
\begin{enumerate}
	\item Scroll to the \textbf{AI Engine API Key} section.
	\item Enter your API key in the password field.
	\item Click \textbf{Save API Key}.
	\item A success notification confirms the update.
\end{enumerate}

% --- Section 9: Status Reference ---------------------------------------------
\section{Status Reference}

The consultation status is shown as a coloured pill on the History page and the Detail page.

\begin{tabular}{m{3cm}m{3cm}m{6cm}}
	\textbf{Status} & \textbf{Colour} & \textbf{Meaning}                                                    \\
	\toprule
	Pending         & Blue-grey       & Audio received; waiting for AI to start.                            \\
	Processing      & Amber           & AI is actively transcribing and analysing.                          \\
	Done            & Teal            & SOAP report is ready.                                               \\
	Failed          & Red             & Processing failed (API error or invalid audio). Retry is available. \\
\end{tabular}

% --- Section 10: Troubleshooting -------------------------------------------
\section{Troubleshooting}

\subsection{I cannot log in}
\begin{itemize}
	\item Verify your email address is typed correctly.
	\item Check that Caps Lock is not enabled.
	\item Contact your administrator to reset your password.
\end{itemize}

\subsection{My audio file fails to upload}
\begin{itemize}
	\item Confirm your file is one of the supported formats: MP3, WAV, M4A, OGG, WebM, FLAC.
	\item Ensure the file is not larger than your organisation's size limit (typically 500 MB).
	\item Try a different browser if the upload consistently fails.
\end{itemize}

\subsection{Processing is stuck on "Pending"}
\begin{itemize}
	\item Wait up to one minute — processing may be queued behind other requests.
	\item If the status remains Pending for more than 5 minutes, refresh the page.
	\item If the problem persists, check with your system administrator.
\end{itemize}

\subsection{Processing failed}
\begin{itemize}
	\item A red \textbf{Failed} card appears on the detail page.
	\item Click \textbf{Retry} to re-process the same audio file — no re-upload needed.
	\item If the retry also fails, note the error message and contact support.
\end{itemize}

\subsection{The SOAP report is in the wrong language}
Use the language tabs (\textbf{Tiếng Việt}, \textbf{English}, \textbf{Français}) above the SOAP sections
to switch the display language. The underlying data is stored in all three languages simultaneously.

% --- Section 11: Glossary ----------------------------------------------------
\section{Glossary}

\begin{tabular}{m{4cm}m{8cm}}
	\textbf{SOAP}            & Acronym for Subjective, Objective, Assessment, Plan — a standard format for clinical documentation. \\
	\textbf{Vina Doctor}     & The web-based portal described in this guide.                                                       \\
	\textbf{AI Scribe}       & The artificial intelligence component that transcribes audio and generates SOAP reports.            \\
	\textbf{DashScope}       & The underlying AI service (Alibaba Cloud) that powers the transcription and analysis.               \\
	\textbf{MP3 / WAV / M4A} & Common audio file format types supported by the upload tool.                                        \\
	\textbf{Qwen}            & The family of large language models used by the AI engine.                                          \\
	\textbf{ICD-10}          & International Classification of Diseases, 10th Revision — diagnostic codes used in the report.      \\
	\textbf{Consultation}    & A single patient encounter, represented by one audio recording and one SOAP report.                 \\
\end{tabular}

\vfill
\color{secondary}\small\raggedright
Vina Doctor AI Scribe — User Guide v1.0 — April 2026 \\
Copyright \textcopyright\ 2026 Vina Doctor. All rights reserved.
\end{document}
